\documentclass[10pt,landscape]{article}

% Versão resumida de naoparametrica.tex: menos assuntos, notação explícita.
% Compilar com XeLaTeX (ou LuaLaTeX): xelatex naoparametrica_resumo.tex
\usepackage{fontspec}
\IfFontExistsTF{Times New Roman}
  {\setmainfont{Times New Roman}}
  {\setmainfont{TeX Gyre Termes}}

\usepackage{multicol}
\usepackage[landscape]{geometry}
\usepackage{amsmath, amssymb}
\usepackage{booktabs}
\usepackage{enumitem}
\setlist[description]{leftmargin=0pt,itemsep=4pt,parsep=0pt,topsep=0pt}
\setlist[itemize]{leftmargin=8pt,itemsep=0pt,parsep=0pt,topsep=1pt}
\usepackage[pdfauthor={Guilherme Marthe},
            pdftitle={Formulário resumido - Estatística Não-Paramétrica}]{hyperref}
\usepackage[open,openlevel=2]{bookmark}

% Notação
\newcommand{\E}{\mathbb{E}}
\newcommand{\V}{\mathbb{V}}
\renewcommand{\P}{\mathbb{P}}
\newcommand{\I}{\mathbb{I}}
\newcommand{\R}{\mathbb{R}}
\newcommand{\IF}{\textrm{IF}}
\newcommand{\VC}{\textrm{VC}}
\newcommand{\MSE}{\textrm{MSE}}
\newcommand{\CV}{\textrm{CV}}
\newcommand{\Fn}{\hat{F}_n}
\newcommand{\X}{\mathbf{X}}
\newcommand{\qc}{\xrightarrow{q.c.}}
\newcommand{\dto}{\xrightarrow{d}}
\newcommand{\pto}{\xrightarrow{P}}

% Geometria e títulos (como em inferencia.tex)
\geometry{top=1cm,left=1cm,right=1cm,bottom=1cm}
\pagestyle{empty}
\makeatletter
\renewcommand{\section}{\@startsection{section}{1}{0mm}%
                                {-1ex plus -.5ex minus -.2ex}%
                                {0.5ex plus .2ex}%
                                {\normalfont\small\bfseries}}
\renewcommand{\subsection}{\@startsection{subsection}{2}{0mm}%
                                {-1explus -.5ex minus -.2ex}%
                                {0.5ex plus .2ex}%
                                {\normalfont\footnotesize\bfseries}}
\makeatother
\setcounter{secnumdepth}{0}
\setlength{\parindent}{0pt}
\setlength{\parskip}{0pt plus 0.5ex}
\newcommand{\secrule}{\smallskip \hrule height 2pt \smallskip}

\begin{document}
\raggedright
\footnotesize
\begin{multicols}{3}
\setlength{\columnsep}{8pt}

{\normalsize\bfseries Formulário resumido -- Estatística Não-Paramétrica}\par
\smallskip
\textbf{Convenção:} $X_1,\dots,X_n$ i.i.d. com função de distribuição (FD) $F$. Toda função aparece com o seu argumento: $F(x)$, $g(x)$, $K(u)$, etc.

%=====================================================================
\section{Ferramentas}\secrule
\begin{description}
  \item[Chebyshev] Se $\E[T]=\theta$: \quad $\P(|T-\theta|\ge t)\le \dfrac{\V(T)}{t^2}$.
  \item[Hoeffding] $Y_1,\dots,Y_n$ independentes com $Y_i\in[0,1]$, $\bar Y=\frac1n\sum_{i=1}^nY_i$:
  \[\P\big(|\bar Y-\E[\bar Y]|\ge\varepsilon\big)\le 2e^{-2n\varepsilon^2}.\]
  \item[Borel-Cantelli] Se $\sum_{n}\P(|T_n-\theta|>\varepsilon)<\infty$ para todo $\varepsilon>0$, então $T_n\qc\theta$. \\
  \emph{Uso:} toda cota do tipo $Ce^{-cn\varepsilon^2}$ é somável em $n$, logo dá convergência quase certa.
  \item[Método delta] Se $\sqrt n(Y_n-\mu)\dto N(0,\sigma^2)$ e $g$ é derivável em $\mu$:
  \[\sqrt n\big(g(Y_n)-g(\mu)\big)\dto N\big(0,\ g'(\mu)^2\sigma^2\big).\]
  \item[Viés e variância] $\MSE(\hat\theta)=\E[(\hat\theta-\theta)^2]=\big(\E[\hat\theta]-\theta\big)^2+\V(\hat\theta)$.
\end{description}

%=====================================================================
\section{Distribuição empírica}\secrule
\begin{description}
  \item[Definição] $\displaystyle \Fn(x)=\frac1n\sum_{i=1}^n\I(X_i\le x)$, \ para cada $x\in\R$.
  \item[Num ponto $x$ fixo] $n\Fn(x)\sim\textrm{Binomial}\big(n,F(x)\big)$, logo
  \[\E[\Fn(x)]=F(x),\qquad \V[\Fn(x)]=\frac{F(x)\big(1-F(x)\big)}{n}\le\frac1{4n}.\]
  \item[Erro máximo] $D_n=\sup_{x\in\R}|\Fn(x)-F(x)|$.
  \item[Glivenko-Cantelli] $D_n\qc0$.
  \item[DKW] $\P(D_n\ge\varepsilon)\le2e^{-2n\varepsilon^2}$, para todo $\varepsilon>0$. \\
  Consequência: quase certamente, $D_n=O\big(\sqrt{\log n/n}\big)$.
\end{description}

\subsection{Molde: consistência por Lipschitz}
Se existe $L>0$ tal que, para quaisquer FDs $G_1,G_2$,
\[|T(G_1)-T(G_2)|\le L\sup_{x}|G_1(x)-G_2(x)|,\]
então $|T(\Fn)-T(F)|\le L\,D_n\qc0$. \\
\emph{Exemplo:} $T(F)=\int_0^MF(x)^2\,dx$. Como $|a^2-b^2|=|a+b||a-b|\le2|a-b|$ para $a,b\in[0,1]$, vale $L=2M$.

\subsection{Molde: mediana amostral}
Se a densidade satisfaz $f(x)\ge C_1>0$, então $|F(x)-F(y)|\ge C_1|x-y|$. Daí
\[|\hat m_n-m|\le\frac{D_n+n^{-1}}{C_1}\qc0.\]

%=====================================================================
\section{Dimensão VC}\secrule
$\mathcal C$ é uma classe de conjuntos e $A=\{x_1,\dots,x_k\}$ é um conjunto finito de pontos.
\begin{description}
  \item[Fragmentar] $\mathcal C$ fragmenta $A$ se \emph{todo} subconjunto de $A$ é da forma $C\cap A$ para algum $C\in\mathcal C$.
  \item[Dimensão VC] $\VC(\mathcal C)=$ o maior $k$ tal que \emph{algum} conjunto de $k$ pontos é fragmentado.
  \item[Coeficiente] $s(\mathcal C,n)=$ número máximo de subconjuntos distintos $C\cap A$, entre conjuntos $A$ com $n$ pontos.
  \item[Sauer-Shelah] Se $\VC(\mathcal C)=d$, então $s(\mathcal C,n)\le(n+1)^d$.
  \item[Desigualdade de VC] Com $P_n(C)=\frac1n\sum_i\I(X_i\in C)$:
  \[\P\Big(\sup_{C\in\mathcal C}|P_n(C)-P(C)|>\varepsilon\Big)\le8\,s(\mathcal C,n)\,e^{-n\varepsilon^2/32}.\]
  Com Sauer-Shelah e Borel-Cantelli: $\VC<\infty\Rightarrow\sup_{C}|P_n(C)-P(C)|\qc0$.
\end{description}
\begin{center}
\begin{tabular}{lc}
\toprule
Classe $\mathcal C$ & $\VC(\mathcal C)$ \\
\midrule
semirretas $(-\infty,a]$ em $\R$ & $1$ \\
intervalos $[a,b]$ em $\R$ & $2$ \\
retângulos $[a_1,b_1]\times[a_2,b_2]$ & $4$ \\
semiespaços em $\R^d$ & $d+1$ \\
\bottomrule
\end{tabular}
\end{center}
\textbf{Receita para mostrar $\VC=d$:} (i) exiba \emph{um} conjunto de $d$ pontos e todos os seus subconjuntos; (ii) para \emph{qualquer} conjunto de $d+1$ pontos, ordene $x_1<\dots<x_{d+1}$ e encontre um subconjunto que nenhum $C$ produz.

%=====================================================================
\section{Funcionais e função de influência}\secrule
\begin{description}
  \item[Funcional] Uma função $T$ que recebe uma FD e devolve um número: $\theta=T(F)$. O estimador plug-in é $T(\Fn)$.
  \item[Funcional linear] Existe uma função $g:\R\to\R$ tal que
  \[T(F)=\int g(x)\,dF(x)=\E_F[g(X)].\]
  O plug-in é a média amostral de $g$: \ $T(\Fn)=\frac1n\sum_{i=1}^ng(X_i)$. \\
  \emph{Exemplos:} média, $g(x)=x$; \ $F(t)$, $g(x)=\I(x\le t)$.
  \item[Massa pontual em $x_0$] É a FD de uma v.a. constante igual a $x_0$:
  \[\delta_{x_0}(y)=\I(y\ge x_0),\qquad y\in\R.\]
  \item[Contaminação] Para $\varepsilon\in[0,1]$, a FD
  \[F_{\varepsilon}(y)=(1-\varepsilon)F(y)+\varepsilon\,\delta_{x_0}(y),\qquad y\in\R,\]
  sorteia de $F$ com prob.\ $1-\varepsilon$ e devolve $x_0$ com prob.\ $\varepsilon$.
  \item[Função de influência] Efeito em $T$ de uma observação extra em $x_0$:
  \[\IF(x_0)=\lim_{\varepsilon\downarrow0}\frac{T(F_\varepsilon)-T(F)}{\varepsilon}=\frac{d}{d\varepsilon}T(F_\varepsilon)\Big|_{\varepsilon=0}.\]
  Sempre vale $\int\IF(x)\,dF(x)=0$: use para conferir a conta.
\end{description}

\begin{center}
\begin{tabular}{ll}
\toprule
$T(F)$ & $\IF(x)$ \\
\midrule
$\int g(y)\,dF(y)$ & $g(x)-T(F)$ \\
média $\mu$ & $x-\mu$ \\
variância $\sigma^2$ & $(x-\mu)^2-\sigma^2$ \\
quantil $\xi_p$ & $\dfrac{p-\I(x\le\xi_p)}{f(\xi_p)}$ \\
\bottomrule
\end{tabular}
\end{center}

\begin{description}
  \item[Regra da cadeia] Se $T(F)=h(\theta_1,\dots,\theta_k)$ com $\theta_j=\int g_j(y)\,dF(y)$:
  \[\IF(x)=\sum_{j=1}^k\frac{\partial h}{\partial\theta_j}(\theta_1,\dots,\theta_k)\,\big(g_j(x)-\theta_j\big).\]
  \item[Normalidade assintótica] Sob regularidade (expansão de von Mises),
  \[\sqrt n\big(T(\Fn)-T(F)\big)\dto N(0,\tau^2),\quad \tau^2=\int\IF(x)^2\,dF(x).\]
  \item[Intervalo de confiança] $\hat\tau^2=\frac1n\sum_{i=1}^n\widehat{\IF}(X_i)^2$ e \
  $T(\Fn)\pm z_{1-\alpha/2}\,\hat\tau/\sqrt n$.
\end{description}

%=====================================================================
\section{Jackknife e bootstrap}\secrule
$\hat\theta_n=T(\Fn)$; \ $\hat\theta_{-i}$ é o mesmo estimador calculado \emph{sem} $X_i$; \ $\bar\theta=\frac1n\sum_{i=1}^n\hat\theta_{-i}$.
\begin{description}
  \item[Jackknife] Viés: \ $\hat b_{\textrm{jack}}=(n-1)(\bar\theta-\hat\theta_n)$. \\
  Estimador corrigido: \ $\hat\theta_{\textrm{jack}}=\hat\theta_n-\hat b_{\textrm{jack}}=n\hat\theta_n-(n-1)\bar\theta$. \\
  Variância: \ $\hat v_{\textrm{jack}}=\frac{n-1}{n}\sum_{i=1}^n(\hat\theta_{-i}-\bar\theta)^2$.
  \item[Por que funciona] Se $\E[\hat\theta_n]=\theta+\frac an+\frac b{n^2}+\dots$, então o viés de $\hat\theta_{\textrm{jack}}$ é $O(n^{-2})$.
  \item[Identidade útil] $\bar X_{-i}=\dfrac{n\bar X-X_i}{n-1}$.
  \item[Bootstrap] Sorteie $X_1^*,\dots,X_n^*$ com reposição de $X_1,\dots,X_n$ e calcule $\hat\theta^*$. Repita $B$ vezes: $\hat\theta^*_1,\dots,\hat\theta^*_B$.
  \item[Erro padrão] $\hat{\textrm{se}}_{\textrm{boot}}^2=\frac1B\sum_{b=1}^B\big(\hat\theta^*_b-\bar\theta^*\big)^2$, \ $\bar\theta^*=\frac1B\sum_b\hat\theta^*_b$.
  \item[Viés] $\hat b^*=\bar\theta^*-\hat\theta_n$.
  \item[Intervalos] normal: $\hat\theta_n\pm z_{1-\alpha/2}\,\hat{\textrm{se}}_{\textrm{boot}}$; \ percentil: quantis $\alpha/2$ e $1-\alpha/2$ dos $\hat\theta^*_b$.
\end{description}

%=====================================================================
\section{Estimador de densidade por kernel}\secrule
\begin{description}
  \item[Kernel] $K:\R\to\R$ com $\int K(u)\,du=1$. Ordem $k$: $\int u^jK(u)\,du=0$ para $j=1,\dots,k$ (todo kernel simétrico tem ordem $1$).
  \item[KDE] Com banda $h>0$:
  \[\hat f_n(x)=\frac1{nh}\sum_{i=1}^nK\Big(\frac{X_i-x}{h}\Big).\]
  \item[Esperança] Trocando $u=(y-x)/h$:
  \[\E[\hat f_n(x)]=\int K(u)\,f(x+uh)\,du.\]
  \item[Suavidade (Hölder)] $f\in\Sigma(\beta,L)$: com $\ell$ o maior inteiro $<\beta$, \ $|f^{(\ell)}(y)-f^{(\ell)}(x)|\le L|y-x|^{\beta-\ell}$. Para $\beta\le1$: $|f(y)-f(x)|\le L|y-x|^\beta$.
  \item[Viés] Kernel de ordem $\ell$ com $\int|u|^\beta|K(u)|\,du<\infty$:
  \[\big|\E[\hat f_n(x)]-f(x)\big|\le C\,h^{\beta}.\]
  \item[Variância] $f$ limitada e $\int K(u)^2du<\infty$: \ $\V[\hat f_n(x)]\le\dfrac{C}{nh}$.
  \item[Banda ótima] $\MSE\le C\big(h^{2\beta}+\frac1{nh}\big)$. Igualando os termos:
  \[h\asymp n^{-1/(2\beta+1)},\qquad \MSE\asymp n^{-2\beta/(2\beta+1)}.\]
  \item[Em $d$ dimensões] variância $\asymp\frac1{nh^d}$, logo $h\asymp n^{-1/(2\beta+d)}$ e $\MSE\asymp n^{-2\beta/(2\beta+d)}$.
  \item[Validação cruzada] $\hat f_{-i}$ é o KDE sem $X_i$:
  \[\CV(h)=\int\hat f_n(x)^2\,dx-\frac2n\sum_{i=1}^n\hat f_{-i}(X_i),\]
  e $\E[\CV(h)]=\E\!\int(\hat f_n(x)-f(x))^2dx-\int f(x)^2dx$.
\end{description}

\subsection{Roteiro para viés/variância/MSE}
\begin{enumerate}[leftmargin=10pt,itemsep=0pt,topsep=1pt]
  \item Escrever a esperança como integral em $y$.
  \item Trocar $u=(y-x)/h$.
  \item Usar $\int K(u)\,du=1$ para subtrair $f(x)$ dentro da integral.
  \item Cotar $|f(x+uh)-f(x)|$ pela suavidade.
  \item Variância: $\V(Z)\le\E[Z^2]$.
  \item Igualar viés$^2$ e variância para achar $h$.
\end{enumerate}

%=====================================================================
\section{Regressão: linear smoothers}\secrule
Modelo: $Y_i=r(X_i)+\varepsilon_i$, \ $\E[\varepsilon_i\mid X_i]=0$, \ $\V[\varepsilon_i\mid X_i]=\sigma^2$.
\begin{description}
  \item[Linear smoother] $\hat r_n(x)=\sum_{i=1}^n\ell_i(x)\,Y_i$, com pesos $\ell_i(x)$ que dependem de $x$ e dos $X_i$, mas não dos $Y_i$.
  \item[Nadaraya-Watson] $\ell_i(x)=\dfrac{K\big((X_i-x)/h\big)}{\sum_{s=1}^nK\big((X_s-x)/h\big)}$.
  \item[Regressograma] $\ell_i(x)=\I(X_i\in B_j)/n_j$ se $x\in B_j$, com $n_j=\#\{i:X_i\in B_j\}$.
  \item[Viés e variância (dado $\X$)] Se $\sum_i\ell_i(x)=1$:
  \[\E[\hat r_n(x)\mid\X]-r(x)=\sum_{i=1}^n\ell_i(x)\big(r(X_i)-r(x)\big),\]
  \[\V[\hat r_n(x)\mid\X]=\sigma^2\sum_{i=1}^n\ell_i(x)^2.\]
  \item[Mesma taxa do KDE] Se os pesos são locais ($\ell_i(x)=0$ quando $|X_i-x|>h$) e $|\ell_i(x)|\le\frac{C}{nh}$: viés $\le Ch^\beta$, variância $\le\frac{C}{nh}$.
  \item[Validação cruzada] Para Nadaraya-Watson, regressograma e MQ, com $L_{ii}=\ell_i(X_i)$:
  \[\CV=\frac1n\sum_{i=1}^n\Big(\frac{Y_i-\hat r_n(X_i)}{1-L_{ii}}\Big)^2.\]
\end{description}

\end{multicols}
\end{document}
